% !TeX root = ../thuthesis-example.tex

\chapter{研究背景与意义}



\section{研究背景}

自机器学习和人工智能流行以来，他们在自然语言处理、图像识别、推荐系统和预测蛋白质结构等方面都取得了惊人的成就 \cite{seq2seq,alphafold}。但与此同时，机器学习模型也变得越来越庞大，模型参数和用来训练的数据集数量也呈现爆炸式增长。 如图\ref{fig:1.1}： 模型大小呈现指数级的上升趋势，在2016年之前，所有模型在18到24个月之间翻倍；而到2016年之后，模型只要3-5个月就会翻倍。面对日益复杂的模型和庞大的数据集，对机器学习模型进行分布式数据并行训练是十分必要的 \cite{edgeCompute}。

\begin{figure}[b]
  \centering
  \includegraphics[width=0.8\linewidth]{figures/1-1.pdf}
  %\caption*{1954年至2021年间流行的新机器学习系统的模型规模。}
  \caption{1954年至2021年间流行的机器学习模型规模（parameters）与发布时间（publication date）变化关系图\cite{modelParameterVsDate}。}
  \label{fig:1.1}
\end{figure}

在分布式训练中，各个节点根据本地的数据集计算局部更新，之后各节点通信，聚合出全局更新累加到全局模型上，从而实现全局模型的更新。为了减少通信数据量以缓解训练过程中的通信瓶颈，分布式训练通过通过稀疏化、量化或低秩的梯度压缩方法，减少了梯度交换的冗余。在传统的独立同分布（IID）分布式环境中，梯度压缩算法可以有效的加速分布式训练，在保证模型精度的条件下实现极高的压缩率，大大缓解了分布式机器学习中的通信瓶颈问题。

但是在新型分布式机器学习环境（如联邦学习 \cite{fed_learning}）中，各节点的数据异质性会让梯度压缩算法出现更严重的精度衰退问题。例如，在IID场景使用Gaia算法（一种应用于跨地域分布式机器学习的梯度压缩算法），设置其超参数为10$\%$，与批量同步并行算法（一种全通信算法，以下简称BSP）相比，精度只降低0.7$\%$。而在非独立同分布（Non-IID）场景中，同样的压缩率设置使精度下降了10.4$\%$。这促使我们设计对节点数据感知的自适应梯度压缩算法，提升梯度压缩算法在Non-IID场景中的鲁棒性。




\section{研究意义与挑战}

有损梯度压缩在Non-IID场景中会造成更严重的精度损失，出现这个现象的原因是他们对不同的节点采用了相同激进的梯度压缩率，而不同的节点往往有不同的数据量和数据分布。例如，在Flickr-mammal（以下简称Flickr）数据集中，拥有大量数据集的节点比拥有第二大数据量的节点多出了78$\%$的数据量（按照次大陆划分节点） \cite{niid2020icml}。在Google Landmark数据集中，拥有最大数据量的节点比其他节点多出至少213$\%$的样本量（按照大洲划分） \cite{Googlelandmark}。

为方便起见，我们将拥有数据量大（小）的节点表示为大（小）节点，将节点拥有的样本数量称为节点尺寸。现存梯度压缩算法中的忽略了节点尺寸的差异，对于那些提出了根据数据分布差异调整压缩率的自适应算法，大节点仍然和小节点使用一样激进的压缩率。在这种情况下，很多可以加速收敛的关键信息就会丢失。

基于我们的测量实验，我们揭示了以下两点对于设计数据感知的梯度压缩策略的看法。首先，在通信受限的Non-IID环境中，各节点相同压缩率的均匀压缩策略并不是最优的，不同尺寸的节点期望分配到不同的压缩率。第二，为了收敛到相同的准确度，给大节点设置更高的压缩率\footnote{压缩率等于压缩后的数据量除以未压缩时的数据量\cite{dc2}}，比起均匀压缩，可以有效节省训练时间。基于这些见解，我们建议设计一种数据感知的自适应梯度压缩算法，以实现在固定且受限通信条件下的最快收敛速率。


\begin{figure}[b]
  \centering
  \includegraphics[width=1.0\linewidth]{figures/1-2.pdf}
  %\caption*{1954年至2021年间流行的新机器学习系统的模型规模。}
  \caption{Flickr-Mammal数据集：每个亚大陆的样本数量\cite{niid2020icml}。}
  \label{fig:1.2}
\end{figure}


%缺少一个flickr/google landmark的图

\iffalse


\section{摘要的写法}

论文摘要是对论文研究内容的高度概括，应具有独立性和自含性，即应是 一篇简短但意义完整的文章。
通过阅读论文摘要，读者应该能够对论文的研究 方法及结论有一个整体性的了解，因此摘要的写法应力求精确简明。
论文摘要 应包括对问题及研究目的的描述、对使用的方法和研究过程进行的简要介绍、 对研究结论的高度凝练等，重点是结果和结论。

论文摘要切忌写成全文的提纲，尤其要避免“第 1 章……；第 2 章……；……”这样的陈述方式。



\section{引言的写法}

一篇学位论文的引言大致包含如下几个部分：
1、问题的提出；
2、选题背 景及意义；
3、文献综述；
4、研究方法；
5、论文结构安排。
\begin{itemize}
  \item 问题的提出：要清晰地阐述所要研究的问题“是什么”。
    \footnote{选题时切记要有“问题意识”，不要选不是问题的问题来研究。}
  \item 选题背景及意义：论述清楚为什么选择这个题目来研究，即阐述该研究对学科发展的贡献、对国计民生的理论与现实意义等。
  \item 文献综述：对本研究主题范围内的文献进行详尽的综合述评，“述”的同时一定要有“评”，指出现有研究状态，仍存在哪些尚待解决的问题，讲出自己的研究有哪些探索性内容。
  \item 研究方法：讲清论文所使用的学术研究方法。
  \item 论文结构安排：介绍本论文的写作结构安排。
\end{itemize}



\section{正文的写法}

本部分是论文作者的研究内容，不能将他人研究成果不加区分地掺和进来。
已经在引言的文献综述部分讲过的内容，这里不需要再重复。
各章之间要存在有机联系，符合逻辑顺序。



\section{结论的写法}

结论是对论文主要研究结果、论点的提炼与概括，应精炼、准确、完整，使读者看后能全面了解论文的意义、目的和工作内容。
结论是最终的、总体的结论，不是正文各章小结的简单重复。
结论应包括论文的核心观点，主要阐述作者的创造性工作及所取得的研究成果在本领域中的地位、作用和意义，交代研究工作的局限，提出未来工作的意见或建议。
同时，要严格区分自己取得的成果与指导教师及他人的学术成果。

在评价自己的研究工作成果时，要实事求是，除非有足够的证据表明自己的研究是“首次”、“领先”、“填补空白”的，否则应避免使用这些或类似词语。

\fi